\clearpage
\begin{examplebox}
\textbf{\textcolor{CExample}{例 0：先看最简单的情况（初中水平）}}
\begin{center}\begin{tikzpicture}[font=\small,>=stealth,line join=round]
\draw[gray!40,rounded corners=3pt] (0,0) rectangle (14.2,8.6);
\node at (7.1,8.0) {例0与快练：两个最小位置相同，可以同时取到};
\begin{scope}[shift={(4.2,4.25)},x=0.8235294117647058cm,y=0.8235294117647058cm]
\draw[->,gray] (-1,0)--(3,0) node[right]{$x$};
\draw[blue,very thick] (0,-2)--(0,3);
\draw[orange!85!black,dashed,thick] (0,-2)--(2,0)--(0,3);
\draw[dotted,gray] (0,0)--(0,-2) (0,0)--(0,3);
\fill (0,-2) circle(2pt) (0,3) circle(2pt);
\fill[blue] (0.0,0) circle(2.3pt);\fill[orange!85!black] (2,0) circle(2pt);
\node[below right=5pt] at (0,-2) {$A$};\node[above right=5pt] at (0,3) {$B$};
\node[below right=8pt] at (0.0,0) {$P_0$};\node[above right=6pt] at (2,0) {$P$};
\end{scope}
\node[align=left,anchor=west] at (8.2,4.9) {$A=(0,-2),\ B=(0,3)$\\[12pt]$P_0=(0,0)$\\[12pt]最短路程 $2+3=5$\\[14pt]一般 $m=n=c$ 时\\[9pt]$x_0=c,\ f_{\min}=a+b$};
\node[text width=12.8cm,align=center] at (7.1,.65) {横纵坐标等比例；蓝线为最短直路，橙色虚线为经其他位置 $P$ 的路。};
\end{tikzpicture}\end{center}



求 $f(x)=\sqrt{x^2+4}+\sqrt{x^2+9}$ 的最小值。

\textbf{分析：} 两个根号都在 $x=0$ 时最小，所以
\[
 f(x)\geq\sqrt4+\sqrt9=2+3=5.
\]
此时两项的最小值刚好可以同时取到，故最小值为 $\boxed5$，对应 $x=0$。

\textbf{为什么还需要 F012？} 当两个根号内分别出现 $(x-m)^2$、$(x-n)^2$，且 $m\ne n$ 时，它们各自最小的位置不同，就不能再把两个最小值直接加起来了。

\medskip
\end{examplebox}
\clearpage
\begin{examplebox}
\textbf{\textcolor{CExample}{例 1：直接识别两个距离（原例保留并补足理解）}}
\begin{center}\begin{tikzpicture}[font=\small,>=stealth,line join=round]
\draw[gray!40,rounded corners=3pt] (0,0) rectangle (14.2,8.6);
\node at (7.1,8.0) {两个根号是两段路：折线拉直，路程最短};
\begin{scope}[shift={(2.5529411764705885,4.25)},x=0.8235294117647058cm,y=0.8235294117647058cm]
\draw[->,gray] (-2,0)--(4,0) node[right]{$x$};
\draw[blue,very thick] (1,-2)--(3,3);
\draw[orange!85!black,dashed,thick] (1,-2)--(-1,0)--(3,3);
\draw[dotted,gray] (1,0)--(1,-2) (3,0)--(3,3);
\fill (1,-2) circle(2pt) (3,3) circle(2pt);
\fill[blue] (1.8,0) circle(2.3pt);\fill[orange!85!black] (-1,0) circle(2pt);
\node[below right=5pt] at (1,-2) {$A$};\node[above right=5pt] at (3,3) {$B$};
\node[below right=8pt] at (1.8,0) {$P_0$};\node[above left=6pt] at (-1,0) {$P$};
\end{scope}
\node[align=left,anchor=west] at (8.2,4.9) {$A=(1,-2)$\\[9pt]$B=(3,3)$\\[12pt]$P_0=(9/5,0)$\\[12pt]$AP_0+P_0B=AB$\\[12pt]最短路程 $\sqrt{29}$};
\node[text width=12.8cm,align=center] at (7.1,.65) {横纵坐标等比例；蓝线为最短直路，橙色虚线为经其他位置 $P$ 的路。};
\end{tikzpicture}\end{center}



求
\[
 f(x)=\sqrt{(x-1)^2+4}+\sqrt{(x-3)^2+9}
\]
的最小值及此时的 $x$。

\textbf{第 1 步（先看结构）：} $m=1$、$n=3$、$a=2$、$b=3$。注意第一个根号单独的最小值 $2$ 在 $x=1$ 取到；第二个根号单独的最小值 $3$ 在 $x=3$ 取到。因此\textbf{不能把 $2+3=5$ 当成原式的最小值}。

\textbf{第 2 步（放置三个点）：} 取 $A(1,-2)$、$B(3,3)$、$P(x,0)$。两点分在轴的两侧，线段 $AB$ 穿过轴，最短时 $AP+PB=AB$。
\[
 f_{\min}=\sqrt{(3-1)^2+(2+3)^2}
 =\boxed{\sqrt{29}}.
\]

\textbf{第 3 步（求交点的位置）：}
\[
 x_0=\frac{3\cdot1+2\cdot3}{2+3}=\boxed{\frac95}.
\]
因为 $1<\dfrac95<3$，此位置恰好在两个定点的横坐标之间。答案：最小值 $\sqrt{29}$，$x=9/5$。

\medskip
\end{examplebox}
\clearpage
\begin{examplebox}
\textbf{\textcolor{CExample}{例 2：隐藏在二次式中的距离（原例保留）}}
\begin{center}\begin{tikzpicture}[font=\small,>=stealth,line join=round]
\draw[gray!40,rounded corners=3pt] (0,0) rectangle (14.2,8.6);
\node at (7.1,8.0) {例2：配方后，定点落在这些真实坐标上};
\begin{scope}[shift={(3.7416666666666667,4.25)},x=0.9166666666666666cm,y=0.9166666666666666cm]
\draw[->,gray] (-2,0)--(4,0) node[right]{$x$};
\draw[blue,very thick] (2,-1)--(-1,2);
\draw[orange!85!black,dashed,thick] (2,-1)--(3,0)--(-1,2);
\draw[dotted,gray] (2,0)--(2,-1) (-1,0)--(-1,2);
\fill (2,-1) circle(2pt) (-1,2) circle(2pt);
\fill[blue] (1.0,0) circle(2.3pt);\fill[orange!85!black] (3,0) circle(2pt);
\node[below right=5pt] at (2,-1) {$A$};\node[above right=5pt] at (-1,2) {$B$};
\node[below left=8pt] at (1.0,0) {$P_0$};\node[above right=6pt] at (3,0) {$P$};
\end{scope}
\node[align=left,anchor=west] at (8.2,4.9) {$A=(2,-1)$\\[9pt]$B=(-1,2)$\\[12pt]$P_0=(1,0)$\\[12pt]$AB=\sqrt{3^2+3^2}$\\[9pt]$f_{\min}=3\sqrt2$};
\node[text width=12.8cm,align=center] at (7.1,.65) {横纵坐标等比例；蓝线为最短直路，橙色虚线为经其他位置 $P$ 的路。};
\end{tikzpicture}\end{center}



求 $f(x)=\sqrt{x^2-4x+5}+\sqrt{x^2+2x+5}$ 的最小值及此时的 $x$。

\textbf{第 1 步（分别配方）：}
\[
 \begin{aligned}
 x^2-4x+5&=(x-2)^2+1^2,\\
 x^2+2x+5&=(x+1)^2+2^2.
 \end{aligned}
\]
因此 $m=2$、$n=-1$、$a=1$、$b=2$。

\textbf{第 2 步（直接算两点距离）：}
\[
 f_{\min}=\sqrt{\bigl(2-(-1)\bigr)^2+(1+2)^2}
 =\sqrt{18}=\boxed{3\sqrt2}.
\]

\textbf{第 3 步（求 $x$ 并代入检查）：}
\[
 x_0=\frac{2\cdot2+1\cdot(-1)}{1+2}=\boxed1.
\]
代回原式：$\sqrt{1-4+5}+\sqrt{1+2+5}=\sqrt2+\sqrt8=3\sqrt2$，正确。

\medskip
\end{examplebox}
\clearpage
\begin{examplebox}
\textbf{\textcolor{CExample}{例 3：已知最小值，反过来求参数（原例保留）}}
\begin{center}\begin{tikzpicture}[font=\small,>=stealth,line join=round]
\draw[gray!40,rounded corners=3pt] (0,0) rectangle (14.2,8.6);
\node at (7.1,8.0) {例3：最短路长为5，横向3与竖向4组成直角三角形};
\begin{scope}[shift={(0.9916666666666671,4.25)},x=0.9166666666666666cm,y=0.9166666666666666cm]
\draw[->,gray] (0,0)--(6,0) node[right]{$x$};
\draw[blue,very thick] (2,-1)--(5,3);
\draw[orange!85!black,dashed,thick] (2,-1)--(1,0)--(5,3);
\draw[dotted,gray] (2,0)--(2,-1) (5,0)--(5,3);
\fill (2,-1) circle(2pt) (5,3) circle(2pt);
\fill[blue] (2.75,0) circle(2.3pt);\fill[orange!85!black] (1,0) circle(2pt);
\node[below right=5pt] at (2,-1) {$A$};\node[above right=5pt] at (5,3) {$B$};
\node[below right=8pt] at (2.75,0) {$P_0$};\node[above left=6pt] at (1,0) {$P$};
\draw[gray,dotted] (2,-1)--(5,-1)--(5,3);\draw (4.8,-1)--(4.8,-.8)--(5,-.8);\node[below=5pt] at (3.5,-1){$3$};\node[right=5pt] at (5,1){$4$};\end{scope}
\node[align=left,anchor=west] at (8.2,4.9) {$A=(2,-1),\ B=(5,3)$\\[12pt]$(5-2)^2+(k+3)^2=25$\\[12pt]$k>0\ \Rightarrow\ k=1$\\[12pt]$P_0=(11/4,0)$};
\node[text width=12.8cm,align=center] at (7.1,.65) {横纵坐标等比例；蓝线为最短直路，橙色虚线为经其他位置 $P$ 的路。};
\end{tikzpicture}\end{center}



已知 $k>0$，
\[
 f(x)=\sqrt{(x-2)^2+k^2}+\sqrt{(x-5)^2+9}
\]
的最小值为 $5$，求 $k$ 及取到最小值时的 $x$。

\textbf{第 1 步（认出参数）：} $m=2$、$n=5$、$a=k$、$b=3$。

\textbf{第 2 步（由最小值列式）：}
\[
 \sqrt{(2-5)^2+(k+3)^2}=5
 \Longrightarrow (k+3)^2=16.
\]
因为 $k>0$，所以 $k+3>0$，只能有 $k+3=4$，故 $\boxed{k=1}$。

\textbf{第 3 步（取等时的 $x$）：}
\[
 x_0=\frac{3\cdot2+1\cdot5}{1+3}=\boxed{\frac{11}{4}}.
\]

\medskip
\end{examplebox}
\clearpage
\begin{examplebox}
\textbf{\textcolor{CExample}{例 4：题目限制 $0\leq x\leq1$，不能照搬公式}}
\begin{center}\begin{tikzpicture}[font=\small,>=stealth,line join=round]
\draw[gray!40,rounded corners=3pt] (0,0) rectangle (14.2,9.5);
\node at (7.1,8.9) {例4：只允许 $0\le x\le1$，最小值换到右端点};
\begin{scope}[shift={(2.0,1.65)},x=1.0cm,y=.72cm]
\draw[->,gray] (-1,0)--(5.6,0) node[right]{$x$};\draw[->,gray] (0,0)--(0,8.4) node[above]{$y$};
\foreach \v/\lab in {2/6,4/8,6/10}{\node[left=4pt] at (0,\v){$\lab$};}\node[below left=5pt] at (0,0){$4$};
\draw[gray,dashed,domain=-1:5.3,samples=150] plot(\x,{sqrt((\x-1)^2+4)+sqrt((\x-3)^2+9)-4});
\draw[blue,very thick,domain=0:1,samples=130] plot(\x,{sqrt((\x-1)^2+4)+sqrt((\x-3)^2+9)-4});
\fill[orange!85!black] (0,{sqrt(5)+sqrt(18)-4}) circle(2.2pt) (1,{2+sqrt(13)-4}) circle(2.2pt);\draw[gray,dotted] (1,0)--(1,{2+sqrt(13)-4});\node[below=5pt] at (1,0){$1$};\end{scope}
\node[align=left,anchor=west] at (8.3,4.9){$0\le x\le1$\\[12pt]原最短位置 $9/5$ 不可用\\[12pt]在这段上一直下降\\[12pt]最大值 $f(0)=\sqrt5+3\sqrt2$\\[10pt]最小值 $f(1)=2+\sqrt{13}$};
\node[text width=12.8cm,align=center] at (7.1,.65) {纵轴从 $4$ 起显示，刻度为原函数值；图取 $m=1,n=3,a=2,b=3$。};
\end{tikzpicture}\end{center}



仍考虑例 1 的函数，但现在要求 $0\leq x\leq1$ 时的值域。

先按完整实数范围求出最佳位置 $x_0=9/5$。由于 $9/5>1$，这次\textbf{根本不允许走到原来的最短位置}。

由证明三可知，函数在 $x<9/5$ 时一直减小。所以在 $[0,1]$ 上，$x=1$ 时最小，$x=0$ 时最大。
\[
 \begin{aligned}
 f(1)&=\sqrt4+\sqrt{4+9}=2+\sqrt{13},\\
 f(0)&=\sqrt{1+4}+\sqrt{9+9}=\sqrt5+3\sqrt2.
 \end{aligned}
\]
因此限制后的\textbf{完整值域}为
\[
 \boxed{\left[2+\sqrt{13},\ \sqrt5+3\sqrt2\right]}.
\]
这个例子说明：\textbf{完整实数范围的最小值公式正确，不代表限制范围后还能取到。}

\medskip
\end{examplebox}
\clearpage
\begin{examplebox}
\textbf{\textcolor{CExample}{例 5：开放的端点，最小值可能根本取不到}}
\begin{center}\begin{tikzpicture}[font=\small,>=stealth,line join=round]
\draw[gray!40,rounded corners=3pt] (0,0) rectangle (14.2,9.5);
\node at (7.1,8.9) {例5：只允许 $x>2$，空心端点不是最小值};
\begin{scope}[shift={(2.0,1.65)},x=1.0cm,y=.72cm]
\draw[->,gray] (-1,0)--(5.6,0) node[right]{$x$};\draw[->,gray] (0,0)--(0,8.4) node[above]{$y$};
\foreach \v/\lab in {2/6,4/8,6/10}{\node[left=4pt] at (0,\v){$\lab$};}\node[below left=5pt] at (0,0){$4$};
\draw[gray,dashed,domain=-1:5.3,samples=150] plot(\x,{sqrt((\x-1)^2+4)+sqrt((\x-3)^2+9)-4});
\draw[blue,very thick,domain=2:5.3,samples=130] plot(\x,{sqrt((\x-1)^2+4)+sqrt((\x-3)^2+9)-4});
\draw[blue,fill=white,thick] (2,{sqrt(5)+sqrt(10)-4}) circle(2.7pt);\draw[gray,dotted] (2,0)--(2,{sqrt(5)+sqrt(10)-4});\node[below=5pt] at (2,0){$2$};\end{scope}
\node[align=left,anchor=west] at (8.3,4.9){$x>2$：端点不允许\\[12pt]蓝线一直上升\\[12pt]只能靠近 $f(2)$\\[12pt]$f(2)=\sqrt5+\sqrt{10}$\\[12pt]没有最小值};
\node[text width=12.8cm,align=center] at (7.1,.65) {纵轴从 $4$ 起显示，刻度为原函数值；图取 $m=1,n=3,a=2,b=3$。};
\end{tikzpicture}\end{center}



仍取例 1 的函数，改为 $x>2$。原来最短位置 $9/5$ 不在范围内，且函数在 $x>9/5$ 时递增。因此 $x$ 越靠近 $2$，函数越小，但因为\textbf{$x=2$ 不允许}，始终不能等于
\[
 f(2)=\sqrt{(2-1)^2+4}+\sqrt{(2-3)^2+9}
 =\sqrt5+\sqrt{10}.
\]
所以这里\textbf{没有最小值}；所有函数值都严格大于 $\sqrt5+\sqrt{10}$，值域为
\[
 \boxed{\left(\sqrt5+\sqrt{10},\ +\infty\right)}.
\]

\medskip
\end{examplebox}
\clearpage
\begin{examplebox}
\textbf{\textcolor{CExample}{快练：对称的特殊情况}}
\begin{center}\begin{tikzpicture}[font=\small,>=stealth,line join=round]
\draw[gray!40,rounded corners=3pt] (0,0) rectangle (14.2,8.6);
\node at (7.1,8.0) {例0与快练：两个最小位置相同，可以同时取到};
\begin{scope}[shift={(4.2,4.25)},x=0.8235294117647058cm,y=0.8235294117647058cm]
\draw[->,gray] (-1,0)--(3,0) node[right]{$x$};
\draw[blue,very thick] (0,-2)--(0,3);
\draw[orange!85!black,dashed,thick] (0,-2)--(2,0)--(0,3);
\draw[dotted,gray] (0,0)--(0,-2) (0,0)--(0,3);
\fill (0,-2) circle(2pt) (0,3) circle(2pt);
\fill[blue] (0.0,0) circle(2.3pt);\fill[orange!85!black] (2,0) circle(2pt);
\node[below right=5pt] at (0,-2) {$A$};\node[above right=5pt] at (0,3) {$B$};
\node[below right=8pt] at (0.0,0) {$P_0$};\node[above right=6pt] at (2,0) {$P$};
\end{scope}
\node[align=left,anchor=west] at (8.2,4.9) {$A=(0,-2),\ B=(0,3)$\\[12pt]$P_0=(0,0)$\\[12pt]最短路程 $2+3=5$\\[14pt]一般 $m=n=c$ 时\\[9pt]$x_0=c,\ f_{\min}=a+b$};
\node[text width=12.8cm,align=center] at (7.1,.65) {横纵坐标等比例；蓝线为最短直路，橙色虚线为经其他位置 $P$ 的路。};
\end{tikzpicture}\end{center}



若 $m=n=c$，则
\[
 f(x)=\sqrt{(x-c)^2+a^2}+\sqrt{(x-c)^2+b^2}\geq a+b,
\]
两项同时在 $x=c$ 时达到各自的最小值，因此 $f_{\min}=a+b$。总公式给出的也是 $\sqrt{0+(a+b)^2}=a+b$，两种看法完全一致。
\end{examplebox}
